%=============================================================================
% MÓDULO 02: RESUMO (Expansão Final)
%=============================================================================
% Frontiers in Public Health — Resumo estruturado (250-400 palavras)
%=============================================================================

\begin{abstract}
\noindent
\textbf{Contexto:} Os modelos de aprendizado de máquina são cada vez mais 
utilizados para o diagnóstico de diabetes, porém o viés algorítmico ameaça 
a prestação equitativa de serviços de saúde. Pesquisas existentes concentraram-se 
predominantemente em atributos demográficos individuais, deixando os efeitos 
interseccionais mal compreendidos.

\textbf{Objetivo:} Este estudo avalia o viés algorítmico em modelos de 
predição de diabetes em múltiplas dimensões demográficas, examinando 
disparidades de desempenho por gênero, etnia e sua intersecção.

\textbf{Métodos:} Analisamos o Banco de Dados de Diabetes dos Índios Pima 
(n=768) utilizando quatro algoritmos de aprendizado de máquina: Regressão 
Logística, Random Forest, Gradient Boosting e Support Vector Machine. Os 
modelos foram avaliados utilizando validação cruzada estratificada em 5 
dobras e conjuntos de teste separados. A equidade foi avaliada utilizando 
Diferença de Odds Equalizados, Diferença de Paridade Demográfica, Disparidade 
de Taxa de Verdadeiros Positivos e Disparidade de Taxa de Falsos Positivos. 
A análise interseccional examinou o desempenho na intersecção de gênero 
e etnia.

\textbf{Resultados:} Todos os quatro modelos apresentaram viés mensurável. 
O modelo SVM alcançou o melhor equilíbrio entre desempenho (F1=0.8313, 
AUC=ROC=0.8794) e equidade. As taxas de falso positivo foram 
consistentemente mais altas para pacientes minoritários (0.333--0.417) 
em comparação com pacientes majoritários (0.164--0.200). A análise 
interseccional revelou viés amplificado para mulheres minoritárias, que 
apresentaram a menor Taxa de Verdadeiros Positivos (0.714) e F1-Score 
(0.750) entre todos os subgrupos. A validação cruzada demonstrou desempenho 
estável entre as dobras (DP < 0.04).

\textbf{Conclusões:} O viés algorítmico em modelos de predição de diabetes 
é sistêmico e multidimensional. Análises de atributo único subestimam a 
extensão da discriminação. Abordagens interseccionais são essenciais para 
identificar desvantagens acumuladas. Recomendamos estratégias multifacetadas 
de mitigação de viés abrangendo coleta de dados, desenvolvimento de 
algoritmos e monitoramento pós-implantação.

\noindent
\textbf{Palavras-chave:} viés algorítmico, diabetes, aprendizado de máquina, 
equidade, interseccionalidade, equidade em saúde, ética em aprendizado 
de máquina, mitigação de viés
\end{abstract}
